% Source: https://artofproblemsolving.com/community/c3h1402934_moab33s_math_tips_4_sum_of_divisors
% Stable identity: aops-topic-1402934-q01
\textbf{(a)} $\sigma(12)=1+2+3+4+6+12=28>24$, so $12$ is abundant. $\sigma(16)=1+2+4+8+16=31<32$, so $16$ is deficient (indeed almost perfect). $\sigma(28)=1+2+4+7+14+28=56=2\cdot28$, so $28$ is perfect.

\textbf{(b)} Since $n$ itself is a divisor of $n$, $\sigma(n)\ge n+1>n$ for every $n\ge2$, so the condition $\sigma(n)>n$ would make every integer from $2$ on ``abundant'' (for example $\sigma(16)=31>16$). The student has mixed up $\sigma(n)$ with the sum of the \emph{proper} divisors, $s(n)=\sigma(n)-n$. The definitions can be written with either function, but consistently: abundant means $s(n)>n$, that is $\sigma(n)>2n$; perfect means $s(n)=n$; almost perfect means $s(n)=n-1$.

\textbf{(c)} The divisors of $2^k$ are $1,2,4,\dots,2^k$, so by the geometric series
\[
\sigma(2^k)=1+2+\dots+2^k=2^{k+1}-1=2\cdot2^k-1 .
\]
Thus every power of $2$ (including $2^0=1$) satisfies $\sigma(n)=2n-1$ and is almost perfect.
